В этой главе сформулирована формальная постановка задачи и обсуждаются используемые метрики качества.

Неформально для того чтобы извлечь библиографию из документа необходимо выполнить следующие шаги:
\begin{enumerate}
    \item Обнаружить в тексте статьи границы библиографической секции. Это подмножество текста публикации назовём $B$.
    \item Из полученного множества текста с библиографической информацией извлекаем строки, которые называются библиографические записи. Множество библиографических записей обозначим, как $R$.
    \item Каждую библиографическую запись разбиваем на поля, называемые библиографические области. В них входят: авторы, цитируемой статьи; ссылаемый заголовок; журнал, где вышла публикация с упомянутой статьёй; год издания номера журнала; страницы в выпуске журнала, на которых изложена цитируемая статья, -- и другие структурные единицы библиографической записи.
\end{enumerate}

Качество каждого этапа оценивается на уровне символов или токенов с помощью $F_1$-меры: насколько точно предсказаны границы каждого поля и записи.

\section{Формализация}

Формально, поставим задачу следующим образом. Пусть задан набор размеченных документов, где $d_i$ -- исходный документ, а именно последовательность токенов с мета-информацией; $y_i$ -- эталонная последовательность меток. Пространство документов $D=\{(d_i,y_i)_{i=1}^m\}$.

Каждый документ $d_i\in D$ представим, как
\begin{equation}
    d_i = \{t_i,x_{i1},m_{i1},\dots,x_{in},m_{in}\},
\end{equation}
где:
\begin{itemize}
    \item Последовательность токенов $t_i$ -- линейный поток слов, чисел, и знаков пунктуации. Он получен путём извлечения из текстового слоя $PDF$-документа; 
    \item Вектор признаков токена $x_{ij}$ -- набор признаков, характеризующий $j-$й токен. В вектор включаются, как геометрические признаки, например координаты-\guillemotleftоффсеты\guillemotright~в текстовом слое, высота шрифта, так и лексико-синтаксические свойства, например, регистр символов, принадлежность к словарю имён или фамилий;
    \item Маска макета $m_{ij}$ -- логический контекст токена, определяющий его положение в структуре страницы. Это может быть индикатор новой строки, признак начала абзаца или попадание в зону колонтитула.
\end{itemize}

Разметка документа $y_i\in Y$ -- последовательность меток длины $n$, вектор целевых классов в формате разметки:
\begin{equation}
    y_i = (y_1^i,\dots,y_n^i),
\end{equation}
где каждая $y_j^i$ принадлежит конечному словарю меток, таким как:
\begin{itemize}
    \item токен вне библиографии,
    \item заголовок в библиографической записи,
    \item автор,
    \item название статьи,
    \item год,
    \item журнал,
    \item страницы.
\end{itemize}

Тогда нужно построить отображение $f:D\to Y$, которое для любого документа предсказывает последовательность меток токенов. Задача заключается в максимизации токен-ориентированной $F_1$-меры.

Короче вышеупомянутые обозначения можно записать так:
$$
\begin{cases}
D &= \{(d_i, y_i)\}_{i=1}^m,\\
d_i &= \{t_i, x_{i1}, m_{i1}, \ldots, x_{in}, m_{in}\},\\
y_i &= (y^i_1, \ldots, y^i_n),\\
f &: D \to Y.
\end{cases}
$$
D — пространство документов; Y — пространство последовательностей меток токенов, максимизирующее символьно-ориентированную $F_1$-меру.

На основе проведенного анализа макро- и микроструктуры документов общая задача семантического анализа и извлечения данных декомпозируется на три последовательных отображения, формирующих сквозной конвейер обработки:

\begin{enumerate}
    \item Макро-локализация области библиографии ($g_1$): Отображение осуществляет бинарную классификацию последовательности токенов документа $t_i$ для выделения непрерывного подмножества $B \subseteq t_i$, содержащего исключительно пристатейный список литературы:
    \[ g_1 : T \to B \]
    
    \item Построчная сегментация библиографических ссылок ($g_2$): Выделенное подмножество $B$ проецируется во множество непересекающихся строк $R = \{r_1, r_2, \dots, r_k\}$, каждая из которых соответствует изолированной библиографической записи:
    \[ g_2 : B \to R \]
    
    \item Семантическое структурирование сущностей ($g_3$): Для каждой извлеченной записи $r \in R$ решается задача распознавания именованных сущностей.. Отображение переводит текстовую строку в валидный структурированный объект $c \in C$ в формате JSON:
    \[ g_3: r \to c \]
\end{enumerate}

Таким образом, результирующее сквозное отображение $f$, преобразующее исходный документ в упорядоченное множество структурированных библиографических записей $C_{\text{doc}} = \{c_1, c_2, \dots, c_k\}$, формально представляется в виде суперпозиции разработанных модулей:
\begin{equation}
    f = g_3 \circ g_2 \circ g_1
\end{equation}
где качество работы каждого промежуточного оператора оптимизируется независимо, а результирующая точность системы оценивается сквозным образом на уровне токенов целевых классов пространства $C$.

Проиллюстрируем введённый формализм на конкретном фрагменте русскоязычной библиографической записи, оформленной в соответствии с ГОСТ.

\textbf{Исходная строка.}
\begin{quote}
\textit{Красоткин С.А. Разработка системы анализа научных публикаций // Системы высокой доступности. --- 2024. --- Т.~15, \textnumero~3. --- С.~45--52.}
\end{quote}

Последовательность токенов $t_i$, извлечённая из текстового слоя, представлена в табл.~\ref{tab:bio-example} совместно с векторами признаков $x_{ij}$ и масками макета $m_{ij}$. Для компактности отображены лексико-синтаксические признаки (принадлежность к словарю фамилий, наличие инициалов, цифровой паттерн) и геометрические оффсеты.

Каждому токену присвоена метка из конечного словаря $\mathcal{Y}$ в формате BIO (Begin-Inside-Outside), что соответствует стандартной схеме кодирования границ именованных сущностей. Метки структурированы следующим образом: префикс \texttt{B} обозначает начало поля, \texttt{I} -- продолжение внутри поля, отсутствие префикса -- токен вне библиографической записи.

\begin{table}[htbp]
\centering
\caption{Пример BIO-разметки библиографической записи}
\label{tab:bio-example}
\small
\begin{tabular}{@{}llp{6.5cm}@{}}
\toprule
\textbf{Токен} & \textbf{Метка $y_j^i$} & \textbf{Пояснение} \\
\midrule
Красоткин & B-AUTH & Начало блока автора (фамилия, словарная принадлежность) \\
С.А. & I-AUTH & Продолжение блока автора (инициалы с точкой) \\
Разработка & B-PAPER & Начало названия статьи (существительное, не в словаре) \\
системы & I-PAPER & Продолжение названия (существительное, омоним функциональному слову) \\
анализа & I-PAPER & Продолжение названия \\
научных & I-PAPER & Продолжение названия \\
публикаций & I-PAPER & Конец названия (граница по двойному слешу \texttt{//}) \\
// & O & Разделитель полей (вне семантических классов) \\
Системы & B-JOUR & Начало названия журнала (заглавная буква после разделителя) \\
высокой & I-JOUR & Продолжение названия журнала \\
доступности & I-JOUR & Конец названия журнала (граница по точке с тире) \\
.--- & O & Пунктуационный разделитель \\
2024 & B-YEAR & Год издания (4-значный цифровой паттерн) \\
.--- & O & Пунктуационный разделитель \\
Т. & B-VOL & Маркер тома (сокращение <<Т.>>) \\
15 & I-VOL & Значение тома (цифровой паттерн) \\
, & O & Пунктуационный разделитель \\
\textnumero & B-ISSUE & Маркер номера выпуска \\
3 & I-ISSUE & Значение номера \\
.--- & O & Пунктуационный разделитель \\
С. & B-PAGES & Маркер страниц (сокращение <<С.>>) \\
45--52 & I-PAGES & Диапазон страниц (цифровой паттерн с тире) \\
\bottomrule
\end{tabular}
\end{table}

Короче, формально постановка задачи записывается следующим образом:
$$
\begin{cases}
g_1 &: t \to B,\\
g_2 &: B \to R,\\
g_3 &: R \to C,\\
f   &= g_3 \circ g_2 \circ g_1
\end{cases}
$$
B — область библиографии; R — множество библиографических записей, C — множество структурированных записей. Оценка проводится по $F_1$ на уровне символов/токенов.

\section{Метрики качества}

Оценка качества композиции отображений $f$ происходит как на макро, так и микро-уровнях за счёт использования симольно-ориентированной $F_1$-меры. Такой способ выбран, потому что стандартная токен-ориентированная метрика не учитывает ошибки в границах слов, вызванные артефактами PDF-конвертации, как отсутствие пробелов или ложные переносы строк.

Пусть $TP_c$ -- количество символов, верно отнесённых к целевому классу, например, заголовок цитируемой статьи; $FP_c$ -- количество символов, ошибочно отнесённых к классу; $FN_c$ -- количество пропущенных символов класса, тогда метрика качества вычисляется следующим образом:
$$
\begin{cases}
F_1 = \frac{2\cdot Precision_c\cdot Recall_c}{Precision_c + Recall_c}, \\
Precision_c=\frac{TP_c}{TP_c+FP_c}, \\
Recall_c = \frac{TP_c}{TP_c + FN_c}.
\end{cases}
$$

Для оценки отображения $g_1$ считается доля правильно найденных библиографическхих блоков. Отображение $g_2$ оценивается по совпадаению границ записей. Считается точность и полнота разбиения на записи. Для оценки отображения $g_3$ используется $F_1$ по каждой библиографической области. Для того чтобы так же учесть частичное перекрытие блоков, возникающее из-за неточных границ полей, используется символьная $F_1$.